\chapterhead{6}{What the Evidence Supports}{Saying what the test showed, no more and no less}\label{ch6}

A fair test gives you evidence. Evidence is not the same as a conclusion. This chapter is about the gap between them.

Here are four sentences someone might write after the paper clip test. They are not equally safe.

\begin{enumerate}[leftmargin=1.6em,itemsep=0.35em]
\item[\textbf{A.}] ``In our test, the clip planes flew 5.3~m on average and the plain planes 4.1~m.''
\item[\textbf{B.}] ``In our test, a clip on the nose made the plane fly about 1.2~m farther.''
\item[\textbf{C.}] ``A paper clip makes every paper plane fly farther.''
\item[\textbf{D.}] ``The clip helps because it moves the weight forward.''
\end{enumerate}

Sentence A is a plain \textbf{report} of the numbers. Sentence B is a careful \textbf{claim}: it stays inside the conditions of the test. Sentence C is too big. We only tested one kind of plane, one clip and one thrower. Sentence D gives a \emph{reason}. Our test never looked at weight, so the reason is a guess, even if it is a good guess.

\begin{center}
\begin{tikzpicture}
  \node[draw=leaf,line width=1.3pt,rounded corners=3pt,fill=leaf!10,text width=2.6cm,align=center,inner sep=5pt,font=\sffamily\small] (a) at (0,0) {\textbf{Report}\\[2pt] what the numbers were};
  \node[draw=sky,line width=1.3pt,rounded corners=3pt,fill=sky!10,text width=2.6cm,align=center,inner sep=5pt,font=\sffamily\small] (b) at (3.5,0) {\textbf{Claim}\\[2pt] what this test showed};
  \node[draw=tangerine,line width=1.3pt,rounded corners=3pt,fill=tangerine!10,text width=2.6cm,align=center,inner sep=5pt,font=\sffamily\small] (c) at (7.0,0) {\textbf{Guess}\\[2pt] why it happened};
  \node[draw=tomato,line width=1.3pt,rounded corners=3pt,fill=tomato!10,text width=2.6cm,align=center,inner sep=5pt,font=\sffamily\small] (d) at (10.5,0) {\textbf{Too big}\\[2pt] always, for every case};
  \draw[-{Stealth},line width=1pt] (a) -- (b); \draw[-{Stealth},line width=1pt] (b) -- (c); \draw[-{Stealth},line width=1pt] (c) -- (d);
\end{tikzpicture}
\end{center}

Think of the ladder from Book~1: observed, counted, estimated, concluded. The further right a sentence is, the more you need to \emph{test it separately}. A report needs only the numbers. A guess about why needs a new test.

\begin{multix}[hmm]\multixbot[hmm]\tcblower
I tried the clip test on one plane, once, and it flew 8~m. So I can now announce: all clips make all planes fly 8~m. Is that science?
\end{multix}

\begin{worked}{An old, famous fair test}
In 1747 a ship's surgeon named James Lind took twelve sailors who all had scurvy and who all ate the same ship's food. He split them into six pairs and gave each pair a different treatment, including cider, vinegar, seawater and oranges and lemons. He reported that the pair who got the oranges and lemons recovered fastest. Lind kept the sailors' food and conditions alike, and changed only the treatment. Twelve people is a very small test by today's standards.
\end{worked}

\subhead{Seven lines for a test report}
The five-line report from Book~1 grows by two lines when a test compares things.

\begin{tcolorbox}[enhanced,colback=white,colframe=ink!60,boxrule=1pt,arc=2mm,left=8pt,right=8pt,top=5pt,bottom=5pt,before skip=7pt,after skip=7pt]
\sffamily\small
\begin{enumerate}[leftmargin=1.6em,itemsep=0.2em]
\item \textbf{Question.} What I wanted to find out.
\item \textbf{What I changed and what I measured.}
\item \textbf{What I kept the same.}
\item \textbf{My readings.} All of them.
\item \textbf{My answer.} Averages, spread and the gap.
\item \textbf{What I claim, and what I do not.}
\item \textbf{What I kept and what it hides.}
\end{enumerate}
\end{tcolorbox}

\clearpage

\begin{worked}{A sample report (practice numbers)}
\textbf{1. Question.} Does a paper clip on the nose make a paper plane fly farther?\\
\textbf{2. Change and measure.} I changed the clip (none, or one small clip). I measured the distance from the launch line to the first landing, nearest 0.1~m.\\
\textbf{3. Kept the same.} The same paper, the same fold, the same thrower, the same launch line and push, the same hall, in a mixed order.\\
\textbf{4. Readings (m).} No clip: 4.1, 3.8, 4.4, 4.0, 4.2. Clip: 5.2, 5.6, 5.1, 5.4, 5.2.\\
\textbf{5. Answer.} No clip about 4.1~m (spread 0.6). Clip about 5.3~m (spread 0.5). Gap 1.2~m.\\
\textbf{6. Claim.} In this test, the clip planes flew about 1.2~m farther. I do not claim this is true for other planes, other clips or other throwers. I do not know why it happened.\\
\textbf{7. Kept and hidden.} I kept the averages, spreads and conditions. I set aside which throw was which, and the exact flight paths. A different thrower or paper may give different numbers.
\end{worked}

\subhead{Which sentence does the evidence support?}
\begin{puzzle}{Safe sentences}
In the ice test, wrapping type, average melting times were: bare 32 min, foil 42, paper towel 57, wool sock 88. Mark each sentence S (supported by the test), G (a guess about why, not tested) or B (too big for the test).
\begin{description}[leftmargin=1.6em,itemsep=0.3em,style=sameline]
\item[(a)] ``In this test, the ice in the sock lasted longest.'' \wline[1cm]
\item[(b)] ``Wool is always the best insulator.'' \wline[1cm]
\item[(c)] ``The sock won because it traps air.'' \wline[1cm]
\item[(d)] ``The bare cube melted in about 32 minutes on average.'' \wline[1cm]
\end{description}
\end{puzzle}

\begin{puzzle}{Repair the claim}
A student writes: ``Flaps make paper planes fly farther.'' The flap test gave flat 4.2~m and flaps 4.3~m, with spreads of 0.9 and 1.1~m.\\
Write a better, more honest sentence.\\
\rule{\linewidth}{0.5pt}\\[4pt]\rule{\linewidth}{0.5pt}
\end{puzzle}

\begin{puzzle}{Find what changed}
A test says: ``We found that plants given plant food grew 3~cm taller. Both plants were beans and both had the same amount of water. One sat in the sun and the other in the shade.'' Is this fair? \wline[2cm] What is the second change? \wline[3.4cm]
\end{puzzle}

\clearpage
\thispagestyle{fancy}
\subhead{Mixed review}
Each puzzle uses something from an earlier chapter.

\begin{puzzle}{Test detective}
Read each story. Find the problem with the test.
\begin{description}[leftmargin=1.6em,itemsep=0.4em,style=sameline]
\item[(a)] ``We wanted to see whether blue paper planes fly farther than red ones. The blue one flew 5~m and the red one 4~m.'' (One throw each.) \wline[3cm]
\item[(b)] ``We tested two kinds of tape for strength. We hung a heavy bag from one and a light bag from the other.'' \wline[3cm]
\item[(c)] ``We threw the plain plane ten times. Then we threw the clip plane ten times. We were tired by the end.'' \wline[3cm]
\end{description}
\end{puzzle}

\begin{puzzle}{Averages again}
Plain plane: 3.8, 4.0, 4.2, 4.0, 4.0~m. Clip plane: 5.1, 5.3, 5.0, 5.4, 5.2~m.\\
(a) Plain: total \wline[1.3cm] average \wline[1.3cm] \quad (b) Clip: total \wline[1.3cm] average \wline[1.3cm]\\
(c) Gap: \wline[1.3cm] \quad (d) Spreads: plain \wline[1.2cm] clip \wline[1.2cm]
\end{puzzle}

\begin{puzzle}{Which claim?}
Two classes ran the ramp test. Class A wrote ``A higher ramp sends cars farther.'' Class B wrote ``In our test, a ramp with two books sent the car about 44~cm farther than one book.'' Which claim is safer? \wline[2cm]\\
Why?\\
\rule{\linewidth}{0.5pt}
\end{puzzle}

\clearpage
\thispagestyle{fancy}\label{slop}
\subhead{Slop check}
Slop is hasty, gaudy work that forgot to take care. It can be made by a person or a machine, quickly or slowly. What matters is whether the work shows care for three things:
\begin{itemize}[leftmargin=1.6em]
\item \textbf{Salient:} Is the main point easy to see?
\item \textbf{Simple:} Is there anything that does nothing but get in the way?
\item \textbf{Relevant:} Does everything on the page help answer the question?
\end{itemize}

Here is a hasty poster about the ice test. Study it, then decide what is wrong.

\begin{center}
\begin{tikzpicture}[scale=0.97]
  \fill[cyan!30,draw=orange!80!black,line width=3pt] (0,0) rectangle (11.4,7.6);
  \node[star,star points=10,fill=magenta!60,draw=blue,line width=2pt,minimum size=2.2cm,font=\bfseries\scriptsize,rotate=-12,text=yellow] at (1.4,6.2) {COOL!!!};
  \node[font=\itshape\bfseries\Large,text=violet,rotate=3] at (6.4,6.8) {IcE iS sO ChIlLy};
  \node[font=\ttfamily\tiny,text=red] at (6.4,6.15) {(no idea what the numbers mean lol)};
  \foreach \x/\h/\c in {2.0/4.4/red,3.4/1.5/green,4.8/3.1/blue,6.2/2.2/orange}{
    \fill[\c,draw=black,line width=0.8pt] (\x,0.9) rectangle (\x+1.0,0.9+\h);
  }
  \fill[yellow,draw=black] (7.8,1) circle (1cm);
  \node[font=\tiny,text=black] at (7.8,1) {a pie?};
  \node[font=\tiny,text=white,rotate=90] at (10.7,3) {ice ice ice ice ice ice ice};
  \foreach \x/\y in {1.2/1.5,1.6/4.2,10.2/5.6,0.8/3}{\node[star,star points=5,fill=yellow,draw=red,minimum size=0.6cm] at (\x,\y) {};}
  \node[font=\tiny,text=lime!50!black] at (5.4,0.45) {look at it!!!  best ever!!!  wow  wow};
\end{tikzpicture}
\end{center}

\begin{slopbox}
\textbf{(a)} What is the poster's main point? Can you tell?\\[3pt]\rule{\linewidth}{0.5pt}\\[6pt]
\textbf{(b)} Name two things that are not simple, because they only distract.\\[3pt]\rule{\linewidth}{0.5pt}\\[6pt]
\textbf{(c)} Name two things that are missing, so that the bars cannot be read.\\[3pt]\rule{\linewidth}{0.5pt}\\[6pt]
\textbf{(d)} Which part is not relevant to the ice test at all?\\[3pt]\rule{\linewidth}{0.5pt}
\end{slopbox}

\clearpage
\thispagestyle{fancy}
\subhead{Now you try}
On a separate sheet, redraw the poster with the same four bars: 32, 42, 57 and 88~minutes. Give it a title that says what was tested, label the bars with the wrapping, label the axis with the unit, pick one or two calm colours, and delete everything that does not help. Check it against the three words: salient, simple, relevant.

Or use the space below for the redraw.

\begin{center}
\begin{tikzpicture}
  \draw[ink!70,line width=1pt,rounded corners=3pt] (0,0) rectangle (11.4,1.0);
  \node[font=\sffamily\small,anchor=west] at (0.2,0.5) {Title:};
  \draw[line width=1.3pt] (1,1.6) -- (1,9.6);\draw[line width=1.3pt] (1,1.6) -- (11,1.6);
  \foreach \y in {2.6,3.6,...,8.6}{\draw[gray!40] (1,\y) -- (11,\y);}
  \node[font=\sffamily\scriptsize,rotate=90] at (0.3,5.6) {label and unit here};
  \node[font=\sffamily\scriptsize] at (6,1.3) {label what the bars stand for};
\end{tikzpicture}
\end{center}

\clearpage
\thispagestyle{fancy}
\subhead{Design your own fair test}
Pick one of these questions, or make up your own: Does a longer ramp send a toy car farther? Does warm water dissolve sugar faster? Does a taller paper tower fall over sooner in a breeze? Fill in the card. Run the test. Then fill in the results table.
\begin{tcolorbox}[enhanced,colback=white,colframe=sky,boxrule=1.2pt,arc=2mm,left=8pt,right=8pt,top=6pt,bottom=6pt,before skip=7pt,after skip=7pt]
\sffamily\small
\textbf{Question.}\\[10pt]\rule{\linewidth}{0.5pt}\\[8pt]
\textbf{I change.}\\[10pt]\rule{\linewidth}{0.5pt}\\[8pt]
\textbf{I measure, with what, to what accuracy.}\\[10pt]\rule{\linewidth}{0.5pt}\\[8pt]
\textbf{I keep the same.}\\[10pt]\rule{\linewidth}{0.5pt}\\[10pt]\rule{\linewidth}{0.5pt}\\[8pt]
\textbf{How many times, and in what order.}\\[10pt]\rule{\linewidth}{0.5pt}\\[4pt]
\end{tcolorbox}
\renewcommand{\arraystretch}{2.2}
\begin{center}\sffamily\scriptsize
\begin{tabular}{|l|c|c|c|c|c|c|}
\hline
 & \textbf{1} & \textbf{2} & \textbf{3} & \textbf{4} & \textbf{5} & \textbf{Average}\\ \hline
\textbf{Version A} & & & & & & \\ \hline
\textbf{Version B} & & & & & & \\ \hline
\end{tabular}
\end{center}
My claim, staying inside the test: \wline[5.5cm]\\[8pt]
\rule{\linewidth}{0.5pt}

\clearpage
\thispagestyle{fancy}
\subhead{What did we keep, across the whole book?}
For each chapter, say one thing that was simplified, what was kept, and what was set aside. A sample is done for you.

\renewcommand{\arraystretch}{2.3}
\begin{center}\sffamily\scriptsize
\begin{tabularx}{\textwidth}{|>{\bfseries}p{2.1cm}|X|X|X|}
\hline
\textbf{Chapter} & \textbf{What we simplified} & \textbf{What we kept} & \textbf{What we set aside} \\ \hline
1 Two planes & two planes thrown in different ways & the question and the one change we wanted to test & the extra differences that hid the answer \\ \hline
2 One change & & & \\ \hline
3 Write it down & & & \\ \hline
4 Again and again & & & \\ \hline
5 Ice in wrappings & & & \\ \hline
6 Evidence & & & \\ \hline
\end{tabularx}
\end{center}

\begin{multix}[wow]\multixbot[wow]\tcblower
Every test you have run in this book was a simplification of something. A fair test keeps most things the same on purpose, so that one difference can show. Next time a headline says something ``works'', you will want to ask: compared with what, and what else was different?
\end{multix}
\clearpage
\thispagestyle{fancy}
\subhead{Tests to try at home}
Here are ten questions you could test with things around the house. Pick three. For each one, write what you would change, what you would measure and one thing you would keep the same.
\begin{enumerate}[leftmargin=1.6em,itemsep=0.2em]
\item Does a ball bounce higher on a hard floor or on a rug?
\item Does warm water or cold water dissolve a spoon of sugar faster?
\item Does a paper cup or a plastic cup keep a hand-warm drink warm longer?
\item Does a longer string make a pendulum swing more slowly?
\item Does a bigger magnet pick up more paper clips?
\item Does a paper towel or a cloth soak up more water?
\item Does a heavier toy car roll farther down the same ramp?
\item Does a plant by a window grow taller than one in a dim corner? (This takes weeks.)
\item Does stirring dissolve salt faster?
\item Do wide or narrow paper strips hold more coins before they sag?
\end{enumerate}
\renewcommand{\arraystretch}{2.4}
\begin{center}\sffamily\scriptsize
\begin{tabularx}{\textwidth}{|p{1.2cm}|X|X|X|}
\hline
\textbf{Number} & \textbf{I change} & \textbf{I measure} & \textbf{I keep the same}\\ \hline
 & & & \\ \hline
 & & & \\ \hline
 & & & \\ \hline
\end{tabularx}
\end{center}
\clearpage
\thispagestyle{fancy}
\subhead{Questions I still have}
Write a question about testing that this book made you wonder about. Then write how you might find out.
\par\vspace{1em}
\foreach \i in {1,...,12}{\noindent\rule{\linewidth}{0.4pt}\par\vspace{20pt}}
\clearpage
